% TOP_PROBLEM: {"problemId": "problem.bass-quillen-conjecture", "canonicalTitle": "Bass-Quillen conjecture", "releaseRank": 312, "primaryDomain": "algebra_representation_category", "primaryDomainLabel": "Algebra, representation & category theory", "displayStatus": "open_with_solved_subcases", "releaseStatus": "open_with_solved_subcases"}
% !TeX program = lualatex
% Definition and sources only; this file is not an LLM solution attempt.
\documentclass[11pt]{article}
\usepackage[margin=25mm]{geometry}
\usepackage{fontspec}
\setmainfont{Latin Modern Roman}
\usepackage{amsmath,amssymb,mathtools}
\usepackage{unicode-math}
\setmathfont{Latin Modern Math}
\usepackage{xurl,hyperref}
\hypersetup{colorlinks=true,urlcolor=blue}
\setcounter{secnumdepth}{0}
\setlength{\parindent}{0pt}
\setlength{\parskip}{5pt}
\setlength{\emergencystretch}{3em}
\begin{document}
\section*{\texorpdfstring{Bass-Quillen conjecture}{Bass-Quillen conjecture}}\label{312}
\subsection{Definitions and mathematical statement}
Let $R$ be a commutative regular Noetherian ring: each localization $R_{\mathfrak p}$ is a regular local ring, whose maximal ideal needs $\dim R_{\mathfrak p}$ generators. A module is finitely generated projective if it is a direct summand of a finite free module. The Bass--Quillen assertion says that for every $d\ge0$ and finitely generated projective $R[t_1,\ldots,t_d]$-module $P$, there exists a finitely generated projective $R$-module $P_0$ with
\[
 P\cong P_0\otimes_R R[t_1,\ldots,t_d].
\]
One may take $P_0=P/(t_1,\ldots,t_d)P$. In geometric language, every vector bundle on affine $d$-space over $R$ is pulled back from $\operatorname{Spec}R$. The Noetherian regular-ring convention is that of the cited formulation; a broader meaning of ``regular'' needs separate specification.
\par\smallskip\textit{Formulation and concept references:} \href{https://arxiv.org/html/2201.06424v4}{[S1]}.

\subsection{Short English statement}
Over a regular ring, must every finite-rank vector bundle on affine space come from a vector bundle on the base ring?
\subsection{Sources}
\begin{itemize}
\item[S1]Kestutis Cesnavicius, Problems about torsors over regular rings with an appendix by Yifei Zhao, arXiv:2201.06424v4, 2025.
\url{https://arxiv.org/html/2201.06424v4}.
\textit{Formulation source.}
\end{itemize}
\end{document}
